% Part: first-order-logic
% Chapter: syntax-and-semantics
% Section: main-operator

\documentclass[../../../include/open-logic-section]{subfiles}

\begin{document}

\olfileid{fol}{syn}{mai}

\olsection{\printtoken{S}{main operator} van een formule}

\begin{explain}
Het is vaak nuttig om te spreken over de laatste operator die bij de
constructie van !!a{formula}~$!A$ is gebruikt. Deze operator heet de
\emph{hoofdoperator} van~$!A$. Intuïtief is dit de ``buitenste''
operator van $!A$. De hoofdoperator van $\lnot !A$ is bijvoorbeeld
$\lnot$, die van $(!A \lor !B)$ is $\lor$, enzovoort.
\end{explain}


\begin{defn}[!!^{main operator}]
\ollabel{def:main-op}
De \emph{!!{main operator}} van !!a{formula}~$!A$ wordt als volgt
gedefinieerd:
\begin{enumerate}
\item \indcase*{!A}{!A}{$\indfrm$ heeft geen !!{main operator}.}

\tagitem{prvNot}{\indcase{!A}{\lnot !B}{de !!{main operator} van $\indfrm$
  is~$\lnot$.}}{}

\tagitem{prvAnd}{\indcase{!A}{(!B \land !C)}{de !!{main operator} van
  $\indfrm$ is~$\land$.}}{}

\tagitem{prvOr}{\indcase{!A}{(!B \lor !C)}{de !!{main operator} van
  $\indfrm$ is~$\lor$.}}{}

\tagitem{prvIf}{\indcase{!A}{(!B \lif !C)}{de !!{main operator} van
  $\indfrm$ is~$\lif$.}}{}

\tagitem{prvIff}{\indcase{!A}{(!B \liff !C)}{de !!{main operator} van
  $\indfrm$ is~$\liff$.}}{}

\tagitem{prvAll}{\indcase{!A}{\lforall[x][!B]}{de !!{main operator}
  van $\indfrm$ is~$\lforall$.}}{}

\tagitem{prvEx}{\indcase{!A}{\lexists[x][!B]}{de !!{main operator} van
  $\indfrm$ is~$\lexists$.}}{}
\end{enumerate}
\end{defn}

In elk geval bedoelen we het specifiek aangegeven \emph{voorkomen} van
de !!{main operator} in de formule. De formule
$((!D \lif !E) \lif (!E \lif !D))$ heeft bijvoorbeeld de vorm $(!B \lif !C)$,
waarbij $!B$ gelijk is aan $(!D \lif !E)$ en $!C$ aan $(!E \lif !D)$.
Het tweede voorkomen van $\lif$ is dus de !!{main operator}.

\begin{explain}
Dit is een \emph{recursieve} definitie van een functie die iedere
niet-atomaire !!{formula} afbeeldt op het voorkomen van haar !!{main
  operator}. Doordat !!{formula}s inductief zijn gedefinieerd, voldoet
iedere !!{formula}~$!A$ aan een van de gevallen in
\olref{def:main-op}. Dit garandeert dat voor iedere niet-atomaire
!!{formula}~$!A$ !!a{main operator} bestaat. Iedere !!{formula} voldoet
aan slechts één van deze voorwaarden en de kleinere !!{formula}s
waaruit $!A$ is geconstrueerd zijn in ieder geval eenduidig bepaald.
Daarom is het voorkomen van de !!{main operator} van~$!A$ uniek en
hebben we dus een functie gedefinieerd.
\end{explain}

We benoemen !!{formula}s zoals in \olref{tab:main-op}, afhankelijk van
het symbool dat hun !!{main operator}
is.\iftag{defNot,defOr,defAnd,defIf,defIff,defTrue,defFalse,defEx,defAll}
{ Merk echter op dat gedefinieerde operatoren officieel niet in
!!{formula}s voorkomen. Het zijn slechts afkortingen en daarom kunnen
zij officieel niet de hoofdoperator van een formule zijn. In bewijzen
over alle !!{formula}s hoeven zij dus niet afzonderlijk te worden
behandeld.}

\begin{table}[!h]
\centering
\begin{tabular}{c | c | c}
!!^{main operator} & Soort !!{formula} & Voorbeeld\\
\hline
geen & atomaire !!{formula} &
\iftag{prvFalse}{$\lfalse$,}{}
\iftag{prvTrue}{$\ltrue$,}{}
$\Atom{R}{t_1, \dots, t_n}$,
$\eq[t_1][t_2]$\\
$\lnot$ & negatie & $\lnot !A$ \\
$\land$ & conjunctie & $(!A \land !B$) \\
$\lor$ & disjunctie & $(!A \lor !B$) \\
$\lif$ & !!{conditional} & $(!A \lif !B$) \\
$\liff$ & !!{biconditional} & $(!A \liff !B)$ \\
$\lforall[][]$ & universele !!{formula}& $\lforall[x][!A]$ \\
$\lexists[][]$ & existentiële !!{formula}& $\lexists[x][!A]$
\end{tabular}
\caption{Hoofdoperator en benamingen van !!{formula}s}
\ollabel{tab:main-op}
\end{table}
\end{document}
